\cleardoubleoddpage
\chapter{Story Summary}
\label{app:story-summary}

This appendix follows the Institute for Machine Learning thesis requirement. Each answer is kept
within three sentences.

\begin{enumerate}
  \item \textbf{What is the central question?}

  Can a compact JAX/Flax latent surrogate learn short-horizon gyrokinetic dynamics and improve
  flux-predictive rollouts over direct persistence of the observed flux on held-out outer-fold
  trajectories?

  \item \textbf{Why is this question important?}

  Nonlinear gyrokinetic simulations are expensive because they evolve a five-dimensional
  distribution function. A compact diagnostic-preserving state could make short temporal studies
  cheaper, but only if it beats a strong and physically meaningful persistence reference.

  \item \textbf{What evidence/data (variables) are needed to answer the question?}

  The evidence is a manifest-/byte-verified universe of 51 trajectories with time-ordered snapshots,
  scalar flux, two one-dimensional spectra, trajectory identifiers, and eight-step rollout windows.
  Split manifests, normalization scope, context, horizon, model-selection rule, and artifact hashes
  are part of the evidence rather than optional metadata.

  \item \textbf{What methods are used to get this evidence/data?}

  A Cyclone/KvikIO adapter validates a channel-first five-dimensional tensor contract, and a
  convolutional encoder with SimSiam and diagnostic objectives produces 128-dimensional latents.
  GRU and causal Transformer sequence models are trained from cached windows, with persistence and
  a frozen diagnostic-head oracle as controls.

  \item \textbf{What analyses must be applied to the data to answer the central question?}

  Five outer group folds and five matched training seeds are evaluated retrospectively, with
  checkpoint selection by validation latent RMSE and family selection by validation flux RMSE only.
  The primary analysis pairs the selected model with observed-flux persistence and uses a
  hierarchical bootstrap over folds, seeds, and trajectories.

  \item \textbf{What evidence/data (values for the variables) were obtained?}

  The accepted ledger contains 230 slots (225 finite, stable metric stages and five selection
  barriers) and 25 explicitly skipped unselected-family test stages. The selected family is Transformer in folds 0, 1, and 3 and GRU
  in folds 2 and 4; every selection record states that test evidence was not opened.

  \item \textbf{What were the results of the analyses?}

  Selected learned flux RMSE is 14.6729 versus 10.0207 for observed persistence, giving a difference
  of +4.6522 with bootstrap interval [+2.5788, +6.6470]. Decoded latent persistence is 14.8049,
  while the diagnostic-head oracle is 12.3732; the learned-minus-latent latent-MSE difference is
  +0.0189 with interval [-0.0699, +0.1332].

  \item \textbf{How did the analyses answer the central question?}

  The learned model does not improve the observed diagnostic baseline in this retrospective nested
  cross-validation estimate. The positive interval and fold-wise positive differences support a
  bounded negative result, not a claim that all latent surrogates fail.

  \item \textbf{What does this answer tell us about the broader field?}

  A complex temporal architecture should not be assumed to beat persistence when the target varies
  slowly and the representation-to-diagnostic path is lossy. The broader lesson is methodological:
  split lineage, validation-only selection, matched seeds, and transparent negative reporting are
  prerequisites for a defensible surrogate claim.
\end{enumerate}
